\chapter{Implications pour le suivi de la fiscalité minière par l'ITIE Sénégal}
\label{chap:recommandations}

L'exercice présenté dans ce rapport, bien que limité à un seul projet et à un seul régime fiscal, met en évidence l'intérêt d'un outil de simulation de type FARI/NRGI pour le travail de l'ITIE Sénégal~: il permet de transformer des paramètres contractuels --- taux de redevance, taux d'impôt, quotités de participation --- en un indicateur de synthèse (\TEMI) directement comparable à celui d'autres juridictions, et de décomposer les recettes attendues par instrument fiscal sur la durée de vie d'un projet. Cette section propose quelques pistes pour l'appropriation de cette démarche par le Comité National ITIE Sénégal, organisées en trois axes~: la consolidation des données d'entrée, le suivi dans le temps, et la transparence sur les arrangements contractuels particuliers.

\section{Consolider les données d'entrée nécessaires à la modélisation}

L'application du modèle à Sabodala-Massawa a mobilisé des données --- calendrier de production, structure des coûts, paramètres fiscaux --- dispersées dans un document technique de plusieurs centaines de pages (le PFS), rédigé en anglais et destiné en premier lieu aux investisseurs et aux marchés financiers. Une première recommandation consiste donc à établir, pour chaque titre minier en phase d'exploitation, une fiche de synthèse normalisée reprenant ces paramètres dans un format directement exploitable par un outil de simulation~: calendrier de production prévisionnel, dépenses d'investissement de développement et de maintien, structure des coûts d'exploitation, et paramètres effectifs du régime fiscal applicable --- en distinguant explicitement les dispositions de droit commun et celles résultant d'une convention particulière. La constitution progressive d'une telle base, pour l'ensemble des titres en exploitation (or, zircon, phosphates, attapulgite), permettrait de reproduire l'exercice réalisé ici pour Sabodala-Massawa à l'échelle du secteur extractif sénégalais dans son ensemble, et de disposer d'une estimation \emph{ex ante} des recettes attendues à comparer aux recettes effectivement déclarées dans les rapports de conciliation.

\section{Assurer un suivi dans le temps du \TEMI{} et de ses déterminants}

Le \TEMI{} calculé dans ce rapport (44,3\,\%) constitue une photographie à un instant donné, fondée sur un prix de l'or de 1\,600~\$US/once et sur le calendrier de production du PFS de 2020. Or, le \TEMI{} d'un projet minier varie avec le prix de la ressource --- principalement via la part relative de l'impôt sur les bénéfices, qui croît plus que proportionnellement avec la profitabilité --- ainsi qu'avec l'arrivée à échéance de certaines exonérations temporaires (TVA, statut d'entreprise franche). Un suivi périodique --- par exemple annuel, à l'occasion de la mise à jour des rapports de conciliation --- de ce \TEMI{}, recalculé sur la base des paramètres de production et de prix effectivement observés, permettrait à l'ITIE Sénégal de documenter l'évolution de la part de la rente minière captée par l'État au fil du cycle de prix des matières premières, et de objectiver le débat --- récurrent dans la plupart des juridictions minières --- sur le caractère « équitable » ou non du partage de cette rente.

\section{Documenter les arrangements contractuels qui échappent au régime fiscal général}

L'écart, souligné au chapitre~\ref{chap:resultats}, entre le \TEMI{} calculé pour le seul régime fiscal général applicable à SGO (44,3\,\%) et l'ordre de grandeur des flux additionnels liés à la transaction Massawa (paiement initial et \emph{gold stream} Franco-Nevada, de l'ordre de 155~millions \$US sur la durée du PFS) illustre une limite plus générale des modèles de type FARI~: ils ne représentent que les instruments fiscaux et parafiscaux « généralisables », et non les arrangements financiers propres à une transaction donnée entre actionnaires privés, qui peuvent pourtant avoir une incidence significative sur la valeur captée localement par rapport à la valeur exportée. La Norme ITIE prévoit déjà la divulgation des contrats et conventions miniers~; ce travail suggère qu'une attention particulière devrait être portée, dans les rapports de conciliation et dans la documentation publiée par l'ITIE Sénégal, à l'identification et, dans la mesure du possible, à la quantification de ces flux contractuels additionnels (contrats de \emph{streaming} ou d'\emph{offtake}, options sur participations, paiements d'étape liés à des acquisitions d'actifs), afin de donner une image complète du partage de la valeur générée par un projet minier, au-delà du seul régime fiscal de droit commun.

\section{Capitaliser sur la onzième colonne ajoutée au modèle}

D'un point de vue pratique, l'ajout d'un onzième régime « Sénégal » au classeur NRGI --- réalisé en complétant une colonne de paramètres déjà prévue par la structure de l'outil, sans modification des formules existantes --- constitue une base directement réutilisable. Il suffit, pour analyser un autre projet minier sénégalais (zircon, phosphates, ou un autre projet aurifère), de répliquer la feuille de paramètres avec le profil opérationnel et les paramètres fiscaux propres à ce projet, le régime « Sénégal » pouvant être affiné --- par exemple en distinguant le régime de droit commun du code minier de 2016 (applicable aux nouveaux titres) du régime conventionnel applicable aux titres antérieurs comme Sabodala. À moyen terme, la constitution d'une bibliothèque de scénarios fiscaux sénégalais --- droit commun 2016, principales conventions minières en vigueur --- au sein d'un même outil permettrait à l'ITIE Sénégal de réaliser, pour tout nouveau projet en cours d'instruction, une comparaison rapide entre le régime envisagé dans la convention et le régime de droit commun, contribuant ainsi à objectiver les négociations entre l'État et les investisseurs évoquées au chapitre~\ref{chap:cadre}.
